\subsection{The statute under the three adversaries}
\label{sec:11.4.4}
Run the statute as \S\ref{sec:11.7} runs the floor:

\begin{itemize}
\item \textbf{An executive willing to carry it out.} The statute prevents: the agency ends and the acts in the schedule stop.
\item \textbf{A captured operator.} The statute binds on paper and raises costs. Funds limitations are law, but the operator can read an account's boundary narrowly, as the NSC staff read Boland's. It can route the act to an agency the statute does not name, and it can spend money appropriated before the statute passed. What survives is the record, and it matters who keeps it. The Iran-Contra committees could not settle the President's role because documents had been shredded and a principal witness had died \autocite{irancontra1987}. In the 2025–26 enforcement operations, by contrast, Amnesty found that video taken by rapid responders and community members proved critical, in several of the most serious incidents, to establishing what happened and challenging official accounts \autocite{Amnesty2026Whistles}. The record was held by neighbours with phones, which is why the schedule protects them.
\item \textbf{Total capture.} The unenacted bill is itself a record of who was asked to end the agency and who signed.
\end{itemize}

The limit is an ordinary statute's. An ordinary statute needs the President's signature or two-thirds of each house, and the next Congress can repeal it. Entrenchment can make repeal a vote on the record, as the entrenched statute of \S\ref{sec:12.1} does; it can't prevent repeal. So the statute's durability lies less in its text than in what it changes before it can be reversed: people released, facilities closed and people returned.
